\subsection{Independent confirmatory energy decision study}
The deployment gate was redesigned using only the preceding Tetouan development case and was then frozen before the PJM test outcomes were inspected. We used four predeclared regional hourly-load series (AEP, COMED, DAYTON, and PJME) from the public Hourly Energy Consumption collection. Duplicate timestamps were averaged, the common period from 1 January 2012 through 2 August 2018 was retained, and each series was aggregated to 2,406 daily means. The data were split chronologically into 20\% prior, 45\% bound, 15\% validation, and 20\% test segments. Scaling, prior localization, and antecedent construction used the prior segment only; certificates used the bound and validation segments; the test segment was excluded from fitting, selection, certification, and deployment decisions.

The frozen candidates were ridge regression, a fixed eight-rule TSK model, and a radius-controlled TSK model, with lags 7, 14, and 28 days. A candidate was deployment-eligible only if its untruncated certificate was at most 0.10, certification target clipping was at most 2\%, validation RMSE and asymmetric cost each improved on a seven-day seasonal-naive fallback by at least 5\%, the improvement was nonnegative in at least three of four chronological validation blocks for both metrics, and the worst block degraded by no more than 5\%. Underforecast errors received weight two and overforecast errors weight one. Among eligible candidates, minimum validation cost determined the decision. All four regional decisions were serialized and hashed before any test prediction was computed. Confidence 0.05 was divided familywise across the four series.
